\part{Bäume}

\chapter{Bäume und Wälder}

\section{Die sparsamste Verbindung}

Satz 2.8 hat gezeigt: Ein zusammenhängender Graph mit $n$ Ecken braucht mindestens $n-1$ Kanten. Und der Weg $P_n$ zeigt, dass $n-1$ auch genügen kann.

Damit stellt sich die Frage nach den Graphen, die diese Schranke genau erreichen: den sparsamsten unter allen, die noch zusammenhängen. Sie heißen Bäume, sie sind der wichtigste Einzelbegriff dieses Bandes, und sie kommen an mehr Stellen vor, als man erwartet.

\begin{merksatz}
\boxtitel{Wo Bäume auftreten}
\begin{center}
{\small
\renewcommand{\arraystretch}{1.45}
\begin{tabular}{@{}P{3.0cm}P{5.4cm}@{}}
\toprule
\textbf{Alarmierungs\-kette} & Jede Person ruft zwei weitere an. Niemand wird doppelt angerufen, niemand vergessen.\\
\textbf{Organigramm} & Jede Stelle hat genau eine vorgesetzte Stelle, die Klinikleitung keine.\\
\textbf{Diagnose\-algorithmus} & Jede Frage verzweigt; von der Eingangsfrage zu jeder Diagnose führt genau ein Weg.\\
\textbf{Stammbaum} & Abstammungsverhältnisse ohne Rückverbindungen.\\
\textbf{Datei\-verzeichnis} & Jede Datei liegt in genau einem Ordner.\\
\bottomrule
\end{tabular}}
\end{center}

\begin{center}
\begin{minipage}{0.88\textwidth}
\kursiv{Das Gemeinsame steht in der dritten Zeile: \enquote{genau ein Weg}. Das ist keine Zufallseigenschaft, sondern eine der fünf gleichwertigen Beschreibungen, die Satz 4.2 aufstellt.}
\end{minipage}
\end{center}
\end{merksatz}

\begin{definition}
\boxtitelm{Definition 4.1 --- Baum, Wald, Blatt}
\begin{itemize}[leftmargin=1.4em, itemsep=4pt]
\item Ein \Begriff{Baum} ist ein Graph, der zusammenhängend und kreisfrei ist.
\item Ein \Begriff{Wald} ist ein Graph, der kreisfrei ist. (Er darf also in mehrere Komponenten zerfallen; jede davon ist dann ein Baum --- daher der Name.)
\item Eine Ecke vom Grad $1$ heißt \Begriff{Blatt}.
\end{itemize}
\end{definition}

\begin{merksatz}
\boxtitel{Alle Bäume mit bis zu fünf Ecken}
\begin{center}
\begin{tikzpicture}[scale=0.9]
\begin{scope}
\node[knoten] (a) at (0,0) {};
\node[font=\scriptsize\itshape, color=grau] at (0,-1.5) {$n=1$};
\end{scope}
\begin{scope}[xshift=1.5cm]
\node[knoten] (b1) at (0,0.4) {};
\node[knoten] (b2) at (0,-0.4) {};
\draw[kante] (b1)--(b2);
\node[font=\scriptsize\itshape, color=grau] at (0,-1.5) {$n=2$};
\end{scope}
\begin{scope}[xshift=3.0cm]
\node[knoten] (c1) at (0,0.5) {};
\node[knoten] (c2) at (0,0) {};
\node[knoten] (c3) at (0,-0.5) {};
\draw[kante] (c1)--(c2); \draw[kante] (c2)--(c3);
\node[font=\scriptsize\itshape, color=grau] at (0,-1.5) {$n=3$};
\end{scope}
\begin{scope}[xshift=4.8cm]
\node[knoten] (d1) at (0,0.75) {};
\node[knoten] (d2) at (0,0.25) {};
\node[knoten] (d3) at (0,-0.25) {};
\node[knoten] (d4) at (0,-0.75) {};
\draw[kante] (d1)--(d2); \draw[kante] (d2)--(d3); \draw[kante] (d3)--(d4);
\node[font=\scriptsize\itshape, color=grau] at (0,-1.5) {$P_4$};
\end{scope}
\begin{scope}[xshift=6.6cm]
\node[knoten] (e0) at (0,0) {};
\node[knoten] (e1) at (0,0.7) {};
\node[knoten] (e2) at (-0.6,-0.5) {};
\node[knoten] (e3) at (0.6,-0.5) {};
\draw[kante] (e0)--(e1); \draw[kante] (e0)--(e2); \draw[kante] (e0)--(e3);
\node[font=\scriptsize\itshape, color=grau] at (0,-1.5) {Stern $K_{1,3}$};
\end{scope}
\begin{scope}[xshift=9.0cm]
\node[knoten] (f1) at (-0.7,0.6) {};
\node[knoten] (f2) at (-0.25,0.2) {};
\node[knoten] (f3) at (0.2,-0.2) {};
\node[knoten] (f4) at (0.65,-0.6) {};
\node[knoten] (f5) at (-0.7,-0.6) {};
\draw[kante] (f1)--(f2); \draw[kante] (f2)--(f3); \draw[kante] (f3)--(f4); \draw[kante] (f2)--(f5);
\node[font=\scriptsize\itshape, color=grau] at (0,-1.5) {einer von dreien mit $n=5$};
\end{scope}
\end{tikzpicture}
\end{center}

Für $n = 1,2,3,4,5$ gibt es genau $1, 1, 1, 2, 3$ verschiedene Bäume (bis auf Isomorphie). Die Zahlen wachsen danach schnell: $6, 11, 23, 47, 106$ für $n = 6, \ldots, 10$. Eine geschlossene Formel dafür ist nicht bekannt --- Abschnitt 5.4 kommt darauf zurück.
\end{merksatz}

\section{Fünf Beschreibungen, die dasselbe bedeuten}

\begin{merksatz}
\boxtitel{Satz 4.2 --- Charakterisierung der Bäume}
Für einen Graphen $T$ mit $n$ Ecken und $m$ Kanten sind die folgenden fünf Aussagen gleichwertig:
\begin{center}
\renewcommand{\arraystretch}{1.55}
\begin{tabular}{@{}P{1.0cm}P{7.4cm}@{}}
\toprule
(i) & $T$ ist ein Baum --- also zusammenhängend und kreisfrei.\\
(ii) & Je zwei Ecken sind durch \kursiv{genau einen} Weg verbunden.\\
(iii) & $T$ ist zusammenhängend und $m = n-1$.\\
(iv) & $T$ ist kreisfrei und $m = n-1$.\\
(v) & $T$ ist zusammenhängend, und jede Kante ist eine Brücke.\\
\bottomrule
\end{tabular}
\end{center}
\end{merksatz}

\begin{merksatz}
\boxtitel{Wie man fünf Aussagen als gleichwertig nachweist}
\enquote{Gleichwertig} heißt: Gilt eine von ihnen, so gelten alle. Streng genommen wären $5 \cdot 4 = 20$ Folgerungen zu zeigen --- von jeder zu jeder anderen.

Es geht sparsamer. Man zeigt einen \Begriff{Ringschluss}:
\[ \text{(i)} \Rightarrow \text{(ii)} \Rightarrow \text{(iii)} \Rightarrow \text{(iv)} \Rightarrow \text{(v)} \Rightarrow \text{(i)} \]
Das sind fünf Beweise statt zwanzig. Warum genügt das? Weil man von jeder Aussage zu jeder anderen gelangt, indem man im Ring weiterläuft --- notfalls einmal ganz herum.

\begin{center}
\begin{tikzpicture}[scale=1.15]
\foreach \i/\lab in {0/{(i)}, 1/{(ii)}, 2/{(iii)}, 3/{(iv)}, 4/{(v)}} {
  \node[knotenb, minimum size=19pt] (r\i) at (90-72*\i:1.25) {\lab};
}
\foreach \i/\j in {0/1,1/2,2/3,3/4,4/0}
  \draw[kante, -{Stealth[length=5pt]}] (r\i) to[bend right=14] (r\j);
\end{tikzpicture}
\end{center}
\end{merksatz}

\begin{merksatz}
\boxtitel{Beweis, Teil 1: (i) $\Rightarrow$ (ii)}
Sei $T$ zusammenhängend und kreisfrei, seien $u,v$ Ecken.

\schritt{Es gibt mindestens einen Weg}
Das ist genau die Definition des Zusammenhangs (Definition 2.4).

\schritt{Es gibt höchstens einen --- Annahme}
Angenommen, es gäbe zwei \kursiv{verschiedene} Wege $P$ und $Q$ von $u$ nach $v$.

\schritt{Die erste Abweichung finden}
Beide beginnen bei $u$. Wir laufen sie gleichzeitig ab. Da sie verschieden sind, gibt es eine erste Stelle, an der sie sich unterscheiden. Sei $x$ die letzte Ecke, bis zu der beide übereinstimmen; danach benutzt $P$ eine Kante $\{x, p\}$ und $Q$ eine Kante $\{x, q\}$ mit $p \neq q$.

\schritt{Die erste Wiedervereinigung finden}
Beide Wege enden bei $v$, treffen sich also wieder. Wir laufen $P$ ab $x$ weiter und schauen, welche Ecke von $P$ als erste wieder auf $Q$ liegt; sei $y$ diese Ecke. Sie existiert, weil spätestens $v$ auf beiden liegt. Und es ist $y \neq x$, denn $y$ liegt hinter $x$ auf $P$ und $P$ hat lauter verschiedene Ecken.

\begin{center}
\begin{tikzpicture}[scale=1.1]
\node[knotenb] (u) at (-2.2,0) {$u$};
\node[knotenb] (x) at (-0.9,0) {$x$};
\node[knotenb] (y) at (1.5,0) {$y$};
\node[knotenb] (v) at (2.8,0) {$v$};
\draw[kante] (u) -- (x);
\draw[kante] (y) -- (v);
\draw[kanted] (x) to[bend left=45] node[above, font=\scriptsize, color=akzent] {$P$} (y);
\draw[kanted] (x) to[bend right=45] node[below, font=\scriptsize, color=akzent] {$Q$} (y);
\node[font=\scriptsize, color=grau] at (-1.55,0.3) {gemeinsam};
\end{tikzpicture}
\end{center}

\schritt{Der Kreis}
Der Abschnitt von $x$ nach $y$ auf $P$ und der Abschnitt von $x$ nach $y$ auf $Q$ haben --- nach Wahl von $x$ und $y$ --- außer $x$ und $y$ selbst keine gemeinsame Ecke. Beide haben Länge $\ge 1$, und mindestens einer hat Länge $\ge 2$: Hätten beide die Länge $1$, so wären es dieselbe Kante $\{x,y\}$, und dann wäre $p = y = q$ im Widerspruch zu $p \neq q$.

Setzt man beide Abschnitte zusammen --- den einen hin, den anderen zurück ---, so entsteht ein geschlossener Kantenzug der Länge $\ge 3$, dessen Ecken bis auf Anfang und Ende verschieden sind. Das ist ein Kreis.

\schritt{Widerspruch}
$T$ war kreisfrei. Also war die Annahme falsch, und es gibt höchstens einen Weg.
\end{merksatz}

\begin{merksatz}
\boxtitel{Beweis, Teil 2: (ii) $\Rightarrow$ (iii)}
\schritt{Zusammenhang}
Aus (ii) folgt sofort: Zwischen je zwei Ecken gibt es einen Weg. Also ist $T$ zusammenhängend.

\schritt{Die Kantenzahl --- Induktion über $n$}
\textbf{Anfang $n=1$:} Ein Graph mit einer Ecke hat keine Kanten, und $n - 1 = 0$. \checkmark

\textbf{Schritt.} Sei $n \ge 2$ und die Behauptung für alle kleineren Eckenzahlen bewiesen. Da $T$ zusammenhängend ist und $n \ge 2$, gibt es mindestens eine Kante; wähle eine, etwa $e = \{x,y\}$.

\schritt{$e$ ist eine Brücke}
Die Kante $e$ selbst ist ein Weg der Länge $1$ von $x$ nach $y$. Nach (ii) ist er der einzige. Läge $e$ auf einem Kreis, so gäbe es nach dem Beweis von Satz 2.14 einen zweiten Weg von $x$ nach $y$ --- den Rest des Kreises. Also liegt $e$ auf keinem Kreis und ist nach Satz 2.14 eine Brücke.

\schritt{$T - e$ hat genau zwei Komponenten}
Mindestens zwei, weil $e$ eine Brücke ist. Höchstens zwei, denn: Sei $w$ eine beliebige Ecke. In $T$ gibt es einen Weg von $w$ nach $x$. Benutzt er $e$ nicht, so liegt $w$ in derselben Komponente von $T-e$ wie $x$. Benutzt er $e$, so muss $e$ der letzte Schritt sein (der Weg endet ja in $x$, und $e$ kommt höchstens einmal vor) --- der Weg läuft also von $w$ nach $y$ ohne $e$ und dann über $e$ nach $x$. Dann liegt $w$ in der Komponente von $y$.

Jede Ecke liegt somit in der Komponente von $x$ oder in der von $y$. Genau zwei.

\schritt{Die Induktionsannahme anwenden}
Seien $T_1$ und $T_2$ diese Komponenten, mit $n_1$ bzw.\ $n_2$ Ecken und $m_1$ bzw.\ $m_2$ Kanten. Es ist $n_1 + n_2 = n$ und $m_1 + m_2 = m - 1$.

Beide erfüllen (ii): Ein Weg zwischen zwei Ecken einer Komponente von $T-e$ ist auch ein Weg in $T$, und ein zweiter dort gäbe einen zweiten in $T$. Beide haben weniger als $n$ Ecken (denn beide sind nichtleer). Nach Induktionsannahme ist $m_1 = n_1 - 1$ und $m_2 = n_2 - 1$.

\schritt{Zusammenrechnen}
\[ m = m_1 + m_2 + 1 = (n_1 - 1) + (n_2 - 1) + 1 = n_1 + n_2 - 1 = n - 1 . \]
\end{merksatz}

\begin{merksatz}
\boxtitel{Beweis, Teil 3: (iii) $\Rightarrow$ (iv)}
Zu zeigen ist nur die Kreisfreiheit; die Kantenzahl ist gegeben.

\schritt{Annahme}
Angenommen, $T$ enthielte einen Kreis $C$. Sei $e$ eine Kante von $C$.

\schritt{Entfernen}
Nach Satz 2.14 ist $e$ keine Brücke, denn $e$ liegt auf einem Kreis. Also hat $T - e$ dieselbe Komponentenzahl wie $T$, ist also weiterhin zusammenhängend.

\schritt{Der Widerspruch}
$T-e$ ist zusammenhängend mit $n$ Ecken und $m - 1 = n - 2$ Kanten. Nach Satz 2.8 braucht ein zusammenhängender Graph mit $n$ Ecken aber mindestens $n-1$ Kanten. Aus $n-2 \ge n-1$ folgt $-2 \ge -1$ --- falsch.

Also gibt es keinen Kreis.
\end{merksatz}

\begin{merksatz}
\boxtitel{Beweis, Teil 4: (iv) $\Rightarrow$ (v)}
\schritt{Schritt 1 --- die Komponenten sind Bäume}
$T$ ist kreisfrei. Jede Zusammenhangskomponente von $T$ ist zusammenhängend (nach Definition) und kreisfrei (ein Kreis in einer Komponente wäre einer in $T$). Also ist jede Komponente ein Baum, erfüllt also (i).

\schritt{Schritt 2 --- ihre Kantenzahlen}
Die Folgerungen (i) $\Rightarrow$ (ii) $\Rightarrow$ (iii) sind bereits bewiesen. Also gilt für jede Komponente $T_i$ mit $n_i$ Ecken und $m_i$ Kanten
\[ m_i = n_i - 1 . \]

\schritt{Schritt 3 --- aufsummieren}
Hat $T$ genau $c$ Komponenten, so ist
\[ m = \sum_{i=1}^{c} m_i = \sum_{i=1}^{c} (n_i - 1) = \Bigl(\sum_{i=1}^{c} n_i\Bigr) - c = n - c . \]

\schritt{Schritt 4 --- $c$ bestimmen}
Nach Voraussetzung ist $m = n-1$. Aus $n - c = n - 1$ folgt $c = 1$. Also ist $T$ zusammenhängend.

\schritt{Schritt 5 --- jede Kante ist Brücke}
$T$ ist kreisfrei, also liegt keine Kante auf einem Kreis. Nach Satz 2.14 ist damit jede Kante eine Brücke.
\end{merksatz}

\begin{merksatz}
\boxtitel{Beweis, Teil 5: (v) $\Rightarrow$ (i)}
$T$ ist nach Voraussetzung zusammenhängend. Zu zeigen bleibt die Kreisfreiheit.

Läge eine Kante $e$ auf einem Kreis, so wäre $e$ nach Satz 2.14 keine Brücke --- im Widerspruch zu (v), wonach \kursiv{jede} Kante eine Brücke ist.

Also gibt es keinen Kreis, und $T$ ist ein Baum. $\blacksquare$

\begin{center}
\begin{minipage}{0.88\textwidth}
\kursiv{Der Ring ist geschlossen. Bemerkenswert an diesem Beweis ist, wie viel Arbeit Satz 2.14 abnimmt: Drei der fünf Teile berufen sich auf ihn. Das ist der Sinn eines Lemmas.}
\end{minipage}
\end{center}
\end{merksatz}

\begin{warnung}
\boxtitelr{Die Bedingung $m = n-1$ allein genügt nicht}
In (iii) und (iv) steht $m = n-1$ jeweils \kursiv{zusammen} mit einer weiteren Bedingung. Lässt man diese weg, wird die Aussage falsch.

\begin{center}
\begin{tikzpicture}[scale=1.1]
\node[knoten] (a) at (-1.6,0.5) {};
\node[knoten] (b) at (-0.9,0.5) {};
\node[knoten] (c) at (-1.25,-0.25) {};
\draw[kante] (a)--(b); \draw[kante] (b)--(c); \draw[kante] (c)--(a);
\node[knoten] (d) at (0.6,0.3) {};
\node[knoten] (e) at (1.4,0.3) {};
\draw[kante] (d)--(e);
\node[knoten] (f) at (2.4,0.3) {};
\node[font=\scriptsize\itshape, color=grau] at (0.4,-1.0) {$n = 6$, $m = 4$? Nein: $m=4$, $n=6$ --- gezählt: $3+1 = 4$ Kanten, $6$ Ecken};
\end{tikzpicture}
\end{center}

Ein sauberes Gegenbeispiel: Nehmen Sie ein Dreieck und eine isolierte Ecke. Dann ist $n = 4$ und $m = 3 = n-1$. Der Graph ist weder zusammenhängend noch kreisfrei --- und damit kein Baum. Die Zahlengleichung stimmt trotzdem.

\begin{center}
\begin{minipage}{0.88\textwidth}
\kursiv{Was hier passiert, sieht man an der Formel $m = n - c$ aus Teil 4: Ein zusätzlicher Kreis erhöht $m$ um eins, eine zusätzliche Komponente senkt es um eins. Beides zusammen hebt sich auf.}
\end{minipage}
\end{center}
\end{warnung}

\section{Blätter}

\begin{merksatz}
\boxtitel{Satz 4.3 --- Jeder Baum mit $n \ge 2$ Ecken hat mindestens zwei Blätter}
\schritt{Beweis --- Schritt 1: einen längsten Weg wählen}
$T$ hat mindestens eine Kante, also gibt es Wege der Länge $\ge 1$. Unter allen Wegen in $T$ wählen wir einen von größtmöglicher Länge:
\[ P = (u_0, u_1, \ldots, u_k), \qquad k \ge 1 . \]
Ein längster existiert, weil ein Weg höchstens $n$ Ecken hat, die Längen also nach oben beschränkt sind.

\schritt{Schritt 2 --- Behauptung: $u_0$ ist ein Blatt}
Angenommen nicht, also $\deg(u_0) \ge 2$. Dann hat $u_0$ außer $u_1$ noch einen weiteren Nachbarn $w$ mit $w \neq u_1$.

\schritt{Schritt 3 --- Fall A: $w$ liegt nicht auf $P$}
Dann ist
\[ (w, u_0, u_1, \ldots, u_k) \]
ein Weg --- alle Ecken sind verschieden, denn $w$ kam auf $P$ nicht vor. Seine Länge ist $k+1 > k$, im Widerspruch zur Wahl von $P$ als längstem Weg.

\schritt{Schritt 4 --- Fall B: $w$ liegt auf $P$}
Dann ist $w = u_j$ für ein $j$. Wegen $w \neq u_1$ und $w \neq u_0$ ist $j \ge 2$. Betrachten wir
\[ (u_0, u_1, \ldots, u_j, u_0) . \]
Das ist ein geschlossener Kantenzug: Die Schritte $u_0 \to u_1 \to \cdots \to u_j$ stammen aus $P$, und der letzte Schritt $u_j \to u_0$ benutzt die Kante $\{u_0, w\} = \{u_0, u_j\}$.

Seine Länge ist $j + 1 \ge 3$, und die Ecken $u_0, \ldots, u_j$ sind paarweise verschieden. Es ist also ein Kreis --- im Widerspruch zur Kreisfreiheit von $T$.

\schritt{Schritt 5 --- Schluss}
Beide Fälle sind unmöglich, also ist $\deg(u_0) = 1$. Dasselbe Argument, angewandt auf das andere Ende, zeigt $\deg(u_k) = 1$. Wegen $k \ge 1$ ist $u_0 \neq u_k$: zwei verschiedene Blätter. $\blacksquare$
\end{merksatz}

\begin{merksatz}
\boxtitel{Folgerung 4.4 --- die mittlere Gradzahl}
In einem Baum mit $n \ge 2$ Ecken ist die Summe aller Grade
\[ \sum_{v} \deg(v) = 2m = 2(n-1) = 2n - 2 . \]
Der Durchschnittsgrad ist also
\[ \frac{2n-2}{n} = 2 - \frac{2}{n} \;<\; 2 . \]

\begin{center}
\begin{minipage}{0.88\textwidth}
\kursiv{Bäume sind knapp unterhalb von Durchschnittsgrad $2$. Da jede Ecke mindestens Grad $1$ hat und der Durchschnitt unter $2$ liegt, muss es Ecken vom Grad $1$ geben --- eine zweite, kürzere Begründung für Satz 4.3, die allerdings nur \kursiv{eine} Ecke liefert. Für \kursiv{zwei} Blätter braucht man das Argument mit dem längsten Weg.}
\end{minipage}
\end{center}
\end{merksatz}

\begin{merksatz}
\boxtitel{Satz 4.5 --- Kante hinzufügen, Kante wegnehmen}
Sei $T$ ein Baum.
\begin{center}
\renewcommand{\arraystretch}{1.55}
\begin{tabular}{@{}P{3.0cm}P{5.6cm}@{}}
\toprule
\textbf{Hinzufügen} & Fügt man zu $T$ eine Kante $\{u,v\} \notin E$ hinzu, so entsteht \kursiv{genau ein} Kreis.\\
\textbf{Wegnehmen} & Entfernt man aus $T$ eine Kante, so zerfällt $T$ in genau zwei Komponenten.\\
\bottomrule
\end{tabular}
\end{center}

\schritt{Beweis, Hinzufügen}
Nach (ii) gibt es in $T$ genau einen Weg $P$ von $u$ nach $v$; seine Länge ist $\ge 2$, weil $\{u,v\}$ keine Kante war. Zusammen mit der neuen Kante ergibt $P$ einen Kreis.

Gäbe es einen zweiten Kreis $C'$ durch die neue Kante, so wäre $C'$ ohne diese Kante ein zweiter Weg von $u$ nach $v$ in $T$ --- Widerspruch zu (ii). Und ein Kreis, der die neue Kante \kursiv{nicht} benutzt, läge ganz in $T$ und widerspräche der Kreisfreiheit.

\schritt{Beweis, Wegnehmen}
Das wurde im Beweis von Teil 2 bereits gezeigt: Jede Kante eines Baumes ist eine Brücke, und $T - e$ hat genau zwei Komponenten. $\blacksquare$
\end{merksatz}

\section{Wurzelbäume}

In den Beispielen des ersten Abschnitts --- Organigramm, Diagnosealgorithmus, Dateiverzeichnis --- gibt es eine ausgezeichnete Ecke: die Klinikleitung, die Eingangsfrage, das oberste Verzeichnis. Das ist eine \kursiv{Zusatzstruktur}, die der Baum selbst nicht mitbringt.

\begin{definition}
\boxtitelm{Definition 4.6 --- Wurzelbaum}
Ein \Begriff{Wurzelbaum} ist ein Paar $(T, r)$ aus einem Baum $T$ und einer ausgezeichneten Ecke $r$, der \Begriff{Wurzel}. Für eine Ecke $v \neq r$ sei $P_v$ der eindeutige Weg von $r$ nach $v$. Dann heißt
\begin{itemize}[leftmargin=1.4em, itemsep=3pt]
\item die vorletzte Ecke auf $P_v$ der \Begriff{Vorgänger} von $v$;
\item $v$ ein \Begriff{Nachfolger} seines Vorgängers;
\item die Länge von $P_v$ die \Begriff{Tiefe} von $v$;
\item die größte vorkommende Tiefe die \Begriff{Höhe} von $(T,r)$.
\end{itemize}
\end{definition}

\begin{merksatz}
\boxtitel{Warum das wohldefiniert ist}
Die Definition beruft sich auf \enquote{den eindeutigen Weg}. Dass es ihn gibt und dass er eindeutig ist, ist genau Aussage (ii) aus Satz 4.2 --- eine Eigenschaft, die wir bewiesen haben und nicht voraussetzen müssen.

\begin{center}
\begin{tikzpicture}[scale=1.1]
\node[knotenr] (r) at (0,1.5) {$r$};
\node[knotenb] (a) at (-1.4,0.6) {};
\node[knotenb] (b) at (0.9,0.6) {};
\node[knotenb] (c) at (-2.1,-0.3) {};
\node[knotenb] (d) at (-0.7,-0.3) {};
\node[knotenb] (e) at (0.9,-0.3) {};
\node[knotenb] (f) at (-2.1,-1.2) {};
\draw[kante] (r)--(a); \draw[kante] (r)--(b);
\draw[kante] (a)--(c); \draw[kante] (a)--(d); \draw[kante] (b)--(e);
\draw[kante] (c)--(f);
\node[font=\scriptsize, color=grau] at (2.5,1.5) {Tiefe $0$};
\node[font=\scriptsize, color=grau] at (2.5,0.6) {Tiefe $1$};
\node[font=\scriptsize, color=grau] at (2.5,-0.3) {Tiefe $2$};
\node[font=\scriptsize, color=grau] at (2.5,-1.2) {Tiefe $3$};
\end{tikzpicture}
\end{center}

Die Höhe dieses Wurzelbaums ist $3$.
\end{merksatz}

\begin{warnung}
\boxtitelr{Wurzel und Höhe sind keine Eigenschaften des Baumes}
Derselbe Baum bekommt je nach Wahl der Wurzel eine andere Höhe.

\begin{center}
\begin{tikzpicture}[scale=1.0]
\begin{scope}
\node[knotenr] (p0) at (0,1.0) {};
\node[knotenb] (p1) at (0,0.35) {};
\node[knotenb] (p2) at (0,-0.3) {};
\node[knotenb] (p3) at (0,-0.95) {};
\node[knotenb] (p4) at (0,-1.6) {};
\draw[kante] (p0)--(p1);\draw[kante] (p1)--(p2);\draw[kante] (p2)--(p3);\draw[kante] (p3)--(p4);
\node[font=\scriptsize\itshape, color=grau] at (0,-2.1) {Wurzel am Ende: Höhe $4$};
\end{scope}
\begin{scope}[xshift=4.6cm]
\node[knotenr] (q0) at (0,0.3) {};
\node[knotenb] (q1) at (-0.75,-0.35) {};
\node[knotenb] (q2) at (0.75,-0.35) {};
\node[knotenb] (q3) at (-0.75,-1.0) {};
\node[knotenb] (q4) at (0.75,-1.0) {};
\draw[kante] (q0)--(q1);\draw[kante] (q1)--(q3);
\draw[kante] (q0)--(q2);\draw[kante] (q2)--(q4);
\node[font=\scriptsize\itshape, color=grau] at (0,-2.1) {Wurzel in der Mitte: Höhe $2$};
\end{scope}
\end{tikzpicture}
\end{center}

Beide Bilder zeigen denselben Baum $P_5$ --- nur die Wurzel ist verschieden gewählt. Aussagen über \enquote{die Höhe eines Baumes} sind daher unvollständig, solange die Wurzel nicht genannt ist.
\end{warnung}

\section{Übungsaufgaben}

\begin{aufgabe}
\boxtitel{Aufgabe 4.1 --- Ein Baum aus seiner Gradfolge}
Ein Baum hat $12$ Ecken. Genau eine Ecke hat Grad $5$, genau zwei haben Grad $3$, und keine Ecke hat Grad $4$ oder mehr als $5$. Wie viele Blätter hat er, und wie viele Ecken vom Grad $2$?

\loesung
\schritt{Schritt 1 --- Die Unbekannten benennen}
Sei $b$ die Anzahl der Blätter (Grad $1$) und $z$ die Anzahl der Ecken vom Grad $2$. Es gibt keine Ecken vom Grad $0$: Ein Baum mit $n \ge 2$ Ecken ist zusammenhängend, also hat jede Ecke mindestens Grad $1$.

\schritt{Schritt 2 --- Erste Gleichung: die Eckenzahl}
\[ b + z + 2 + 1 = 12 \quad\Longrightarrow\quad b + z = 9 \]
(die $2$ steht für die beiden Ecken vom Grad $3$, die $1$ für die vom Grad $5$).

\schritt{Schritt 3 --- Zweite Gleichung: die Gradsumme}
Nach Folgerung 4.4 ist die Gradsumme $2n - 2 = 2 \cdot 12 - 2 = 22$. Also
\[ 1 \cdot b + 2 \cdot z + 3 \cdot 2 + 5 \cdot 1 = 22 \quad\Longrightarrow\quad b + 2z + 11 = 22 \quad\Longrightarrow\quad b + 2z = 11 . \]

\schritt{Schritt 4 --- Das Gleichungssystem lösen}
Wir ziehen die erste Gleichung von der zweiten ab:
\begin{align*}
(b + 2z) - (b + z) &= 11 - 9\\
z &= 2
\end{align*}
Einsetzen in $b + z = 9$ ergibt $b = 7$.

\schritt{Antwort}
$\mathbf{7}$ Blätter und $\mathbf{2}$ Ecken vom Grad $2$.

\schritt{Probe}
Eckenzahl: $7 + 2 + 2 + 1 = 12$ \checkmark\\
Gradsumme: $7\cdot 1 + 2\cdot 2 + 2\cdot 3 + 1 \cdot 5 = 7 + 4 + 6 + 5 = 22 = 2 \cdot 11 = 2(n-1)$ \checkmark

\schritt{Ein solcher Baum, hingezeichnet}
\begin{center}
\begin{tikzpicture}[scale=1.05]
\node[knotenb] (z) at (0,0) {$5$};
\node[knotenb] (a) at (-1.5,0.9) {$3$};
\node[knotenb] (b) at (1.5,0.9) {$3$};
\node[knoten] (l1) at (-0.7,-1.1) {};
\node[knoten] (l2) at (0.0,-1.2) {};
\node[knoten] (l3) at (0.7,-1.1) {};
\node[knotenb] (m1) at (-2.5,0.3) {$2$};
\node[knoten] (l4) at (-3.3,-0.4) {};
\node[knoten] (l5) at (-2.1,1.8) {};
\node[knotenb] (m2) at (2.5,0.3) {$2$};
\node[knoten] (l6) at (3.3,-0.4) {};
\node[knoten] (l7) at (2.1,1.8) {};
\draw[kante] (z)--(a); \draw[kante] (z)--(b);
\draw[kante] (z)--(l1); \draw[kante] (z)--(l2); \draw[kante] (z)--(l3);
\draw[kante] (a)--(m1); \draw[kante] (a)--(l5); \draw[kante] (m1)--(l4);
\draw[kante] (b)--(m2); \draw[kante] (b)--(l7); \draw[kante] (m2)--(l6);
\end{tikzpicture}
\end{center}
Nachzählen: $11$ Kanten, $12$ Ecken. \checkmark
\end{aufgabe}

\begin{aufgabe}
\boxtitel{Aufgabe 4.2 --- Eine Ungleichung für Wurzelbäume}
Sei $(T,r)$ ein Wurzelbaum der Höhe $h$, in dem jede Ecke höchstens $k$ Nachfolger hat. Zeigen Sie
\[ n \;\le\; 1 + k + k^2 + \cdots + k^h . \]
Welche Höhe braucht ein solcher Baum mindestens, wenn $n = 1000$ und $k = 2$?

\loesung
\schritt{Schritt 1 --- Nach Tiefen sortieren}
Jede Ecke hat genau eine Tiefe. Sei $a_t$ die Anzahl der Ecken der Tiefe $t$. Dann ist
\[ n = a_0 + a_1 + \cdots + a_h . \]

\schritt{Schritt 2 --- $a_0 = 1$}
Nur die Wurzel hat Tiefe $0$, denn für $v \neq r$ hat der Weg von $r$ nach $v$ Länge $\ge 1$.

\schritt{Schritt 3 --- Der Zusammenhang zwischen den Ebenen}
Jede Ecke $v$ der Tiefe $t+1$ hat einen Vorgänger. Behauptung: Dieser hat Tiefe $t$.

Begründung: Sei $P_v = (r, \ldots, u, v)$ der eindeutige Weg zu $v$, mit Vorgänger $u$. Dann ist $(r, \ldots, u)$ ein Weg von $r$ nach $u$ --- und wegen der Eindeutigkeit \kursiv{der} Weg. Seine Länge ist um $1$ kleiner als die von $P_v$, also hat $u$ die Tiefe $t$.

\schritt{Schritt 4 --- Abschätzen}
Jede Ecke der Tiefe $t+1$ ist Nachfolger einer Ecke der Tiefe $t$, und verschiedene Ecken der Tiefe $t+1$ haben verschiedene\ldots{} nein: Sie können denselben Vorgänger haben. Aber jede Ecke der Tiefe $t$ hat höchstens $k$ Nachfolger. Also
\[ a_{t+1} \;\le\; k \cdot a_t . \]

\schritt{Schritt 5 --- Auflösen}
Aus $a_0 = 1$ folgt der Reihe nach $a_1 \le k$, $a_2 \le k \cdot a_1 \le k^2$, allgemein $a_t \le k^t$. Also
\[ n = \sum_{t=0}^{h} a_t \;\le\; \sum_{t=0}^{h} k^t = 1 + k + \cdots + k^h . \]

\schritt{Schritt 6 --- Der Zahlenwert}
Für $k = 2$ ist die geometrische Summe (Band 3, Kapitel 2)
\[ 1 + 2 + 4 + \cdots + 2^h = 2^{h+1} - 1 . \]
Gefordert ist $2^{h+1} - 1 \ge 1000$, also $2^{h+1} \ge 1001$.

\begin{center}
\renewcommand{\arraystretch}{1.4}
\begin{tabular}{@{}P{1.6cm}P{2.4cm}P{2.6cm}@{}}
\toprule
$h$ & $2^{h+1}-1$ & \textbf{reicht?}\\
\midrule
$8$ & $511$ & nein\\
$\mathbf{9}$ & $\mathbf{1023}$ & \textcolor{gruen}{ja}\\
\bottomrule
\end{tabular}
\end{center}

\schritt{Antwort}
Mindestens Höhe $\mathbf{9}$.

\begin{center}
\begin{minipage}{0.86\textwidth}
\kursiv{Das ist die Rechnung hinter der Bemerkung, dass man in einem Suchbaum mit tausend Einträgen etwa zehn Fragen braucht. Sie ist zugleich der Grund, warum Diagnosealgorithmen mit Ja-Nein-Fragen erstaunlich flach bleiben: Verdoppelt sich die Zahl der Diagnosen, wächst die Höhe nur um eins.}
\end{minipage}
\end{center}
\end{aufgabe}

\begin{aufgabe}
\boxtitel{Aufgabe 4.3 --- Wälder}
Ein Wald hat $n$ Ecken, $m$ Kanten und $c$ Komponenten. Zeigen Sie $m = n - c$, und bestimmen Sie die Anzahl der Komponenten eines Waldes mit $20$ Ecken und $14$ Kanten.

\loesung
\schritt{Schritt 1 --- Die Komponenten sind Bäume}
Jede Komponente eines Waldes ist zusammenhängend (das ist die Definition der Komponente) und kreisfrei (ein Kreis darin wäre einer im Wald). Also ist jede ein Baum.

\schritt{Schritt 2 --- Kantenzahl je Komponente}
Nach Satz 4.2, Aussage (iii), hat die $i$-te Komponente $m_i = n_i - 1$ Kanten.

\schritt{Schritt 3 --- Summieren}
Da jede Kante in genau einer Komponente liegt:
\[ m = \sum_{i=1}^{c} m_i = \sum_{i=1}^{c}(n_i - 1) = n - c . \]

\schritt{Schritt 4 --- Der Zahlenwert}
$14 = 20 - c$, also $c = \mathbf{6}$ Komponenten.

\schritt{Probe an einem Beispiel}
Sechs Bäume mit zusammen $20$ Ecken: etwa fünf Wege $P_3$ ($5 \cdot 3 = 15$ Ecken, $5 \cdot 2 = 10$ Kanten) und ein $P_5$ ($5$ Ecken, $4$ Kanten). Zusammen $20$ Ecken, $14$ Kanten, $6$ Komponenten. \checkmark
\end{aufgabe}

\begin{aufgabe}
\boxtitel{Aufgabe 4.4 --- Zwei Bäume mit gleicher Gradfolge}
Geben Sie zwei nicht isomorphe Bäume mit derselben Gradfolge an. Was folgt daraus für die Invariantenliste aus Kapitel~1?

\loesung
\schritt{Schritt 1 --- Wonach man sucht}
Gesucht sind zwei Bäume mit gleicher Eckenzahl (dann automatisch gleicher Kantenzahl, denn $m = n-1$) und gleicher Gradfolge, die sich dennoch unterscheiden.

\schritt{Schritt 2 --- Die kleinsten Beispiele durchprobieren}
Für $n \le 5$ ist die Gradfolge eines Baumes bereits eindeutig bestimmend --- man kann alle Fälle durchgehen. Bei $n = 6$ wird man fündig.

Wir suchen eine Gradfolge mit Summe $2n-2 = 10$ bei sechs Einträgen. Nehmen wir $(3,2,2,1,1,1)$: Summe $= 3+2+2+1+1+1 = 10$. \checkmark

\schritt{Schritt 3 --- Zwei Bäume dazu}
\begin{center}
\begin{tikzpicture}[scale=1.1]
\begin{scope}
\node[knotenb] (a) at (0,0) {$3$};
\node[knotenb] (b) at (-1.0,0.7) {$2$};
\node[knotenb] (c) at (1.0,0.7) {$2$};
\node[knoten] (d) at (0,-0.9) {};
\node[knoten] (e) at (-1.8,1.3) {};
\node[knoten] (f) at (1.8,1.3) {};
\draw[kante] (a)--(b); \draw[kante] (a)--(c); \draw[kante] (a)--(d);
\draw[kante] (b)--(e); \draw[kante] (c)--(f);
\node[font=\scriptsize\itshape, color=grau] at (0,-1.5) {$T_1$: die Zweien getrennt};
\end{scope}
\begin{scope}[xshift=5.6cm]
\node[knotenb] (p) at (0,0) {$3$};
\node[knotenb] (q) at (1.0,0.7) {$2$};
\node[knotenb] (s) at (1.8,1.4) {$2$};
\node[knoten] (t) at (2.6,2.1) {};
\node[knoten] (u) at (-0.9,0.6) {};
\node[knoten] (w) at (-0.9,-0.6) {};
\draw[kante] (p)--(q); \draw[kante] (q)--(s); \draw[kante] (s)--(t);
\draw[kante] (p)--(u); \draw[kante] (p)--(w);
\node[font=\scriptsize\itshape, color=grau] at (0.8,-1.5) {$T_2$: die Zweien hintereinander};
\end{scope}
\end{tikzpicture}
\end{center}

\schritt{Schritt 4 --- Nachweis, dass sie nicht isomorph sind}
Wir brauchen eine Invariante, die sie unterscheidet.

\kursiv{Erster Versuch --- die Nachbarn der Ecke vom Grad 3.} In $T_1$ hat die Ecke vom Grad $3$ die Nachbargrade $2, 2, 1$. In $T_2$ hat sie die Nachbargrade $2, 1, 1$. Die Liste der Nachbargrade ist eine Invariante, denn ein Isomorphismus erhält Grade und Nachbarschaften.

Also $T_1 \not\cong T_2$.

\kursiv{Zweiter Versuch zur Kontrolle --- der Durchmesser.} In $T_1$ ist der längste Weg $e - b - a - c - f$, Länge $4$. In $T_2$ ist der längste Weg $u - p - q - s - t$, Länge $4$. Diese Invariante unterscheidet \kursiv{nicht} --- ein gutes Beispiel dafür, dass man mehrere braucht.

\schritt{Schritt 5 --- Was daraus folgt}
Die Gradfolge steht in der Invariantenliste aus Kapitel~1 an dritter Stelle. Sie ist also auch bei Bäumen nicht ausreichend, obwohl Bäume die einfachsten zusammenhängenden Graphen sind.

\begin{center}
\begin{minipage}{0.86\textwidth}
\kursiv{Dass sogar Bäume nicht durch ihre Gradfolge bestimmt sind, macht deutlich, wie weit man von einem praktikablen Isomorphie\-kriterium entfernt ist. Für Bäume gibt es allerdings --- anders als für Graphen allgemein --- ein schnelles Verfahren, das Isomorphie sicher entscheidet; Kapitel~15 nennt es.}
\end{minipage}
\end{center}
\end{aufgabe}

\begin{geschichte}
\boxtitelg{Cayley, die Chemie und die Bäume}
\Person{Arthur Cayley} (1821--1895) führte den Begriff des Baumes 1857 ein, um Differentialoperatoren zu ordnen. Zwanzig Jahre später kam er auf ein ganz anderes Problem zurück: die Zahl der \kursiv{Isomere} gesättigter Kohlenwasserstoffe.

Ein Alkan $\mathrm{C}_k\mathrm{H}_{2k+2}$ lässt sich als Graph auffassen: Kohlenstoffatome sind Ecken vom Grad $4$, Wasserstoffatome Ecken vom Grad $1$. Lässt man die Wasserstoffatome weg, bleibt ein Baum, in dem jede Ecke höchstens vier Nachbarn hat. Die Zahl der Isomere ist die Zahl solcher Bäume bis auf Isomorphie.

Cayley zählte sie 1875 für $k$ bis $13$. Seine Zahlen stimmten mit den damals bekannten Isomeren überein --- ein bemerkenswerter Fall, in dem reine Kombinatorik eine chemische Vorhersage machte.

\begin{center}
{\small
\renewcommand{\arraystretch}{1.4}
\begin{tabular}{@{}P{2.2cm}P{1.6cm}P{1.6cm}P{1.6cm}@{}}
\toprule
\textbf{$k$ Kohlenstoffe} & $4$ & $5$ & $6$\\
\midrule
\textbf{Isomere} & $2$ & $3$ & $5$\\
\bottomrule
\end{tabular}}
\end{center}

\begin{center}
\begin{minipage}{0.9\textwidth}
\kursiv{Die Zählung unbewurzelter Bäume ist schwer, weil man Symmetrien herausrechnen muss. Cayley machte dabei für $k = 13$ einen Fehler; er wurde erst 1931 bemerkt. Das Verfahren, das solche Zählungen zuverlässig macht, stammt von \Person{George Pólya} (1937) und steht in Band 15, Kapitel~9.}
\end{minipage}
\end{center}
\end{geschichte}

\begin{merksatz}
\boxtitel{Was Sie jetzt können}
\begin{center}
\renewcommand{\arraystretch}{1.5}
\begin{tabular}{@{}P{3.0cm}P{5.6cm}@{}}
\toprule
\textbf{Satz 4.2} & Fünf gleichwertige Beschreibungen eines Baumes, im Ringschluss bewiesen\\
\textbf{Satz 4.3} & Jeder Baum mit $n \ge 2$ hat mindestens zwei Blätter\\
\textbf{Satz 4.5} & Kante hinzufügen erzeugt genau einen Kreis; Kante wegnehmen trennt in zwei Teile\\
\textbf{Formel} & $m = n - c$ für jeden Wald\\
\textbf{Wurzelbaum} & Tiefe, Höhe --- und dass beides von der Wurzelwahl abhängt\\
\bottomrule
\end{tabular}
\end{center}

\begin{center}
\begin{minipage}{0.88\textwidth}
\kursiv{Im nächsten Kapitel geht es um Bäume \kursiv{innerhalb} eines Graphen: Wie findet man einen sparsamsten Teilgraphen, der noch alle Ecken verbindet --- und wie viele davon gibt es?}
\end{minipage}
\end{center}
\end{merksatz}
